\subsection{Dimension of Simple Representations}

Denote the cardinal of the group G by n.~Let \(\pi\) be a linear representation of G in a finite-dimensional K-vector space M. For every \(g \in G\) , we have \(\pi(g)^{n} = 1_{M}\) , so the minimal polynomial of \(\pi(g)\) divides \(T^{n} - 1\) . Since \(n \cdot 1 \neq 0\) in K, this minimal polynomial is separable (V, §11, No.~2, p.~78), and since the field K is algebraically closed, the endomorphism \(\pi(g)\) of M is diagonalizable (VII, §5, No.~7, p.~40, Proposition 12). The eigenvalues of \(\pi(g)\) are n-th roots of unity, and for every \(\alpha \in K\) , the geometric multiplicity of \(\alpha\) as an eigenvalue of \(\pi(g)\) (VII, §5, No.~2, p.~30, Definition 1) is equal to the multiplicity of \(\alpha\) as a root of the characteristic polynomial of \(\pi(g)\) . Denote by \(\mathcal{O}_{n}\) the subgroup of K generated by the set \(\mu_{n}(K)\) of n-th roots of unity; it is a finitely generated Z-module and a subring of K. The character of \(\pi\) takes its values in \(\mathcal{O}_{n}\) .

PROPOSITION 9. --- Suppose that the field K has characteristic zero. Then the degree of every simple representation of G divides the cardinal of G.

Let \((\mathrm{V},\pi)\) be a simple representation of G and \(\chi\) be its character. For every element a of \(\mathrm{Z}(\mathrm{K}[\mathrm{G}])\) , the endomorphism \(\pi(a)\) of V is a homothety (VIII, p.~47, Theorem 1); denote by \(\varphi(a)\) the scalar such that \(\pi(a)=\varphi(a)_{\mathrm{V}}\) .

The resulting mapping \(\varphi\) from \(Z(K[G])\) to \(K\) is an algebra homomorphism. Let us take \(a = \sum_{g \in G} \chi(g^{-1})g\) ; by the remark of VIII, p.~408, we have \(\varphi(a) = (\dim V)^{-1}|G|\) . On the other hand, \(a\) belongs to the subring \(\mathcal{O}_n[G] \cap Z(K[G])\) of \(K[G]\) , which is a finitely generated \(\mathbf{Z}\) -module (VII, §3, p.~15, Corollary). So the element \(\varphi(a) = (\dim V)^{-1}|G|\) of \(K\) belongs to a subring of \(K\) that is a finitely generated \(\mathbf{Z}\) -module. We conclude using the following lemma.

Lemma. --- Let L be an extension of Q. Let A be a subring of L. Suppose that A is a finitely generated Z-module. We have \(A \cap Q = Z\) .

Since the Z-module \(A \cap Q\) is finitely generated, there exists a strictly positive integer N such that \(A \cap Q\) is contained in \(\frac{1}{N}Z\) . Let x be an element of \(Q - Z\) ; we write it as \(x = \frac{p}{q}\) , where p and q are mutually prime integers and \(q \geqslant 2\) . We have \(q^{N} \geqslant 2^{N} > N\) (Set Theory, III, §3, No.~6, p.~165, Theorem 2), the integers \(p^{N}\) and \(q^{N}\) are mutually prime, and consequently \(x^{N} \notin \frac{1}{N}Z\) . It follows that x does not belong to A. This concludes the proof of the lemma.

In the next subsection, we extend Proposition 9 to the case when we only assume that the characteristic of K does not divide the order of G.

